% ECONOMICS_PROBLEM: {"id": "D73-385", "title": "Durable anticorruption policy", "jel_code": "D73", "source_id": "OP-0383"}
% FIRST_POSTED: 2026-10-09
% ATTEMPT_STATUS: partial
% ENTRY_KIND: research_attempt
\documentclass[11pt]{article}
\usepackage[T1]{fontenc}
\usepackage[utf8]{inputenc}
\usepackage[margin=25mm]{geometry}
\usepackage{amsmath,amssymb,amsthm,xurl,hyperref}
\newtheorem{proposition}{Proposition}
\newtheorem{lemma}{Lemma}
\newtheorem{corollary}{Corollary}
\newcommand{\E}{\mathbb E}
\newcommand{\R}{\mathbb R}
\newcommand{\Prb}{\mathbb P}
\newcommand{\Var}{\operatorname{Var}}
\newcommand{\Cov}{\operatorname{Cov}}
\DeclareMathOperator*{\argmax}{arg\,max}
\DeclareMathOperator*{\argmin}{arg\,min}
\setlength{\parindent}{0pt}
\setlength{\parskip}{6pt}
\title{Durable anticorruption policy: a research attempt}
\author{GPT-6 (Codex)}
\date{9 October 2026}
\begin{document}
\maketitle

\section*{Target and outcome}
\textbf{Status: partial; durable harm reductions for an actual organization remain unresolved.}
The target is a finite-horizon change in true, weighted corruption harm, accompanied by service and enforcement outcomes. I derive a sharp deterrence threshold in a strategic avenue-choice model, an explicit harmful-substitution example, detection-adjusted bounds, and a dynamic adaptation formula. Together they show why neither fewer targeted offenses nor more detected offenses determines the sign of the requested effect.

Olken and Pande's primary manuscript explicitly discusses adaptation and substitution following anticorruption policies \cite{op}. The parameters and propositions below are additional analytical specializations. No claim is made that a particular government's policy has one of the calculated effects.

\section*{Officials choose the least deterred avenue}
An official can act honestly or choose one of $K$ mutually exclusive corrupt acts. Honest payoff is zero. Avenue $k$ produces private rent $r_k>0$, sanction $s_k>0$ conditional on discovery, and detection probability $d_k\in[0,1]$. The official is risk neutral, knows the policy, and chooses the act with highest expected payoff
\[
 u_k=r_k-s_kd_k.
\]
Ties with honesty are resolved in favor of honesty; ties between corrupt avenues use a declared fixed rule. True harm from avenue $k$ is $\lambda_k\geq0$, independent of the fine, which is a transfer under this accounting. Audit capacity is $\sum_kd_k\leq B$. Sanctions are fixed legal inputs rather than unbounded controls.

\begin{proposition}[Feasibility of complete deterrence]
A feasible policy deters every avenue exactly when
\[
 r_k/s_k\leq1\quad\hbox{for every }k,
 \qquad\sum_kr_k/s_k\leq B.
\]
The smallest total audit probability that achieves weak deterrence is $\sum_kr_k/s_k$.
\end{proposition}
\begin{proof}
Honesty is selected exactly when $u_k\leq0$ for every avenue, requiring $d_k\geq r_k/s_k$. Probability and capacity constraints yield necessity. Setting every $d_k$ at its threshold establishes sufficiency and the minimum sum. If honesty is not selected at indifference, strict deterrence requires strict inequalities and spare capacity; the displayed value is then an infimum.
\end{proof}

For uncertain rents $r_k\in[\underline r_k,\overline r_k]$ and known sanctions, the same argument gives robust deterrence thresholds $\overline r_k/s_k$. This is a minimax feasibility result with an explicit uncertainty set. It is not necessarily welfare-optimal: complete deterrence can cost more than the harm it prevents, or reduce service if audits consume production resources.

\section*{A targeted success with higher total harm}
Take two avenues with rents $(4,3)$, sanctions $(10,10)$, and harm weights $(1,3)$. Without audits the official chooses avenue 1, producing harm 1. A policy auditing only avenue 1 with $d_1=0.2$ changes its net payoff to 2, so the official switches to avenue 2 with payoff 3 and harm 3. Targeted corruption disappears, while total harm increases by 2. The numerical example remains strict under small perturbations, so it does not depend on an indifference tie.

A capacity $B=0.7$ permits weak complete deterrence using $(d_1,d_2)=(0.4,0.3)$ under the specified honesty tie rule. A smaller capacity cannot deter both in this model. A weighted welfare objective could nonetheless prefer a smaller audit bundle if enforcement costs are high. Measuring the substitution avenue and the service consequences is therefore necessary even when the targeted outcome is perfectly observed.

\section*{Discovery is a separate observation process}
For a fixed avenue and policy regime, let true binary incidence be $C$, detected incidence $D$, sensitivity $s=\Prb(D=1\mid C=1)$, and false-positive rate $f=\Prb(D=1\mid C=0)$. These $s,f$ are measurement parameters, distinct from the sanction notation above. If $c=\Prb(C=1)$ and $d=\Prb(D=1)$,
\[
 d=f+(s-f)c.
\]
With known $s>f$, true incidence is $(d-f)/(s-f)$. Thus changing audits changes the mapping from incidence to measured offenses. For instance, true incidence can fall from $0.4$ to $0.2$ while sensitivity rises from $0.2$ to $0.8$, yielding detected rates $0.08$ and $0.16$ when false positives vanish. Detected offenses double while true incidence halves.

Suppose independent validation instead supplies intervals
\[
 f\in[f_L,f_U],\qquad s\in[s_L,s_U],\qquad
 0\leq f_L\leq f_U<s_L\leq s_U\leq1.
\]
The observed rate must satisfy $d\in[f_L,s_U]$ for compatibility. The sharp incidence interval in this measurement class is
\[
 c_L=\max\left\{0,\frac{d-f_U}{s_U-f_U}\right\},\qquad
 c_U=\min\left\{1,\frac{d-f_L}{s_L-f_L}\right\}.
\]
To see the lower endpoint, choose maximal sensitivity and maximal feasible false positives. If $d\leq f_U$, choosing $f=d$ gives $c=0$; otherwise use $(f_U,s_U)$. For the upper endpoint choose minimal false positives and minimal feasible sensitivity. If $d\geq s_L$, choose $s=d$ to obtain $c=1$; otherwise use $(f_L,s_L)$. Continuity fills the interval. This proof also checks endpoint attainability rather than merely reporting an outer bound.

If a policy is randomized across independent organizational clusters and validation is representative within each arm, the arm-specific detected laws and validation restrictions identify such intervals for true arm incidence. For nonnegative fixed harm weights, a valid effect interval is
\[
 \left[\sum_k\lambda_k(c^1_{k,L}-c^0_{k,U}),
       \sum_k\lambda_k(c^1_{k,U}-c^0_{k,L})\right].
\]
It is sharp when avenue-specific opportunities and measurement mechanisms can vary independently subject to these marginal restrictions. If one imposes the earlier one-act choice model, incidence probabilities must also sum to at most one; then joint optimization under those restrictions can tighten the interval. It would be incorrect to claim marginal sharpness automatically proves joint sharpness. False positives, validation selection, and audit spillovers require corresponding additional restrictions.

\section*{A finite-state adaptation benchmark}
Let $x_t$ be expected counts in finitely many active corruption states, with transitions
\[
 x_{t+1}=A_px_t+b_p,
\]
where $A_p$ is a nonnegative column-substochastic matrix (each column sums to at most one) and $b_p\geq0$ represents entry into corrupt states. Exits include deterrence, dismissal, and retirement. Assume $\rho(A_p)<1$. Then
\[
 x_T=A_p^Tx_0+\sum_{s=0}^{T-1}A_p^sb_p,
 \quad
 x_\infty=(I-A_p)^{-1}b_p.
\]
Multiplying the regime difference by the harm-weight vector gives the exact horizon-specific effect. A short-run decline from $A_px_0$ does not determine the long-run difference because policy can change both substitution transitions and entry. If a policy merely moves officials into a persistent, unmeasured state, the observed initial effect can be opposite to the complete-system effect. The stability assumption is substantive; if an absorbing corrupt state receives continuing entry, the stationary formula does not exist.

Repeated cohort data with policy variation could identify transition and entry changes, but detected-state transitions are not automatically true-state transitions. Independent audits or a separately identified hidden-state model are needed. Officials moving to uncovered agencies must either remain in the target system or be explicitly excluded from the estimand; changing the boundary after observing displacement changes the policy question.

\section*{Remaining task}
The mathematical examples leave the organization's real rents, sanction enforceability, costs, response transitions, and validation quality unknown. A full assessment must track all material avenues, service output, personnel selection, and policy learning over the chosen horizon. The results here establish restrictions under which those data could be interpreted and give exact counterexamples to two common shortcuts; they do not identify the optimal anticorruption bundle without those data.

\section{Completion Estimate}
58\%

This is a post hoc, heuristic estimate for the full catalogue target.
Strategic avenue choice gives a deterrence threshold and harmful-substitution example, with sharp detection-adjusted incidence bounds and horizon-specific adaptation formulas. Actual rents, enforceable sanctions, personnel transitions, service outcomes and interagency displacement remain unidentified for the durable policy comparison.

\begin{thebibliography}{9}
\bibitem{op} B. A. Olken and R. Pande (2012), \emph{Corruption in Developing Countries}, February manuscript, abstract and discussion of adaptation. Primary MIT-hosted PDF inspected on 9 October 2026. \url{https://economics.mit.edu/research/publications/corruption-developing-countries}.
\end{thebibliography}
\end{document}
